\subsection{Consistency in the weak class}\label{sec:B5T05}
\begin{teorema}[Mesoscopic consistency]\label{thm:meso}
For fixed $C,\eta,\gamma$ there exist $\varepsilon_0,k_0,C_*$ independent
of $n/k$ with the following property. If the parent satisfies the conditions of Section~\ref{sec:B5T01} with
$\varepsilon\le\varepsilon_0$ and $k\ge k_0$, then all its legal children
have the small stationary point constructed in Section~\ref{sec:B5T03}, and
\[
 \boxed{\left|\log\rho(B)\right|\le C_*/k^2,\qquad
 \rho(B)=\sum_{c\in r}\exp[L(B+c)-L(B)].}
\]
\end{teorema}
The statement fixes a row $r$ and includes all its children. Here $L$ is precisely
the canonical predictor, with its tree term, local factor, RC cofactor, and covolume.
We do not identify its corrected marginal with the true marginal of the solutions.
The predictor selects the small branch. The negative Hessian
is positive definite in the convex region $p<1/4$ with positive vacancies.
Thus the Newton point is its unique stationary point on that branch and requires no new selection.
The total displacement is $y=z+s+O_{\rm rel}(k^{-2})$, where $s$ is the Newton step from $\bar p$.
A margin is reserved in $\varepsilon_0,C,\eta,\gamma$ for all children.

\begin{proof}
\emph{Adaptation of the macroscopic proof.}
The proof follows that of Section~\ref{sec:macro} and retains its algebraic identities and normalization.
In place of its uses of $k\ge\delta_0n$, we use Section~\ref{sec:B5T02} for the line inverse
and the response, Section~\ref{sec:B5T03} for the Newton step, the mass of the Schur difference,
and the cubic remainders, and Section~\ref{sec:B5T04} for the traces, their changes, and the resolvent.
Rebalancing uses only the $2k$ RC lines.
At no point do we require the number of diagonals to be $O(k)$. The correspondence is as follows.
\begin{center}
\renewcommand{\arraystretch}{1.25}
\begin{tabular}{@{}>{\raggedright\arraybackslash}p{0.52\textwidth}
                >{\raggedright\arraybackslash}p{0.43\textwidth}@{}}
\toprule
Step of the macroscopic proof & Replaced by\\
\midrule
$G=I+O_{\max}(k^{-1})$ and $\|G\|_\infty=O(1)$ (Section~\ref{sec:B3T05}) &
 Section~\ref{sec:B5T02}, Schur complement\\
$\|P^{-1}R\|_\infty=O(1)$ (Section~\ref{sec:B3T06}) &
 Section~\ref{sec:B5T02}, weighted response\\
Newton step and cubic remainder of the action (Section~\ref{sec:B4T02}) &
 Section~\ref{sec:B5T03}\\
$\|\mathcal E\|_F^2$, $\sum_{ij}|\mathcal E_{ij}|$ and $\operatorname{tr}W_c$
 (Sections~\ref{sec:B4T01} and~\ref{sec:B4T04}) &
 Section~\ref{sec:B5T03}, full Schur difference\\
Changes of $\sum p_i^3$, $\sum p_i^2\kappa_i$ and $\sum_am_{2,a}^2$ (Section~\ref{sec:B4T07}) &
 Section~\ref{sec:B5T04}, prefactor and entrywise resolvent\\
Quadratic term and cubic remainder of the determinant (Section~\ref{sec:B4T08}) &
 Section~\ref{sec:B5T04}, change of support\\
\bottomrule
\end{tabular}
\end{center}

\par\emph{Exact change of support and gauge.}
Let $\Gamma=S_0^+$, $v=(\mathbf1_R,-\mathbf1_C,0_D)$,
$\Pi=vv^T/(2k)$, and $\widetilde S_0=S_0+\Pi$.
Let $J$ be the four removed lines and let $I=J^c$. Set
\[
 C_J=S_{0,JJ}+\Pi_{JJ},\quad V_J=S_{0,IJ}+\Pi_{IJ},
 \quad P_*=(\widetilde S_0)_{II}-V_JC_J^{-1}V_J^T,
\]
\[
 H_*=P_*^{-1}=(\Gamma+\Pi)_{II},\quad
 D_*=\Delta_0+(\Pi'-\Pi_{II})+V_JC_J^{-1}V_J^T.
\]
The exact cofactor identity is
\[
 \frac{\det S'_{0,\rm gf}}{\det S_{0,\rm gf}}
 =\frac{k}{k-1}\frac{\det(I+H_*D_*)}{\det C_J}.
\]
Section~\ref{sec:B5T04} controls its quadratic and cubic terms. No Taylor expansion is taken
all the way to $p=0$ or $1$.

\par\emph{Preservation of the cancellation and its coefficients.}
Let $a_i$ be the indicator of the four lines of $i$. Define
$f_i=\sum_{a<b\in i}\Gamma_{ab}$ and, using only the retained pairs, $f_i^I$.
The quadratic response of the child satisfies
\[
 y_i/p_i=-a_i^T\Gamma a_c+O(k^{-2}),\qquad
 \mathcal Q(y)=a_c^T\Gamma a_c-4+12p_c+O(k^{-2}).
\]
Using the exact Jacobi derivative avoids differentiating an asymptotic remainder and gives
\[
 \delta_c=-6p_c-f_c+\sum_Fp_i(6p_i+f_i)
             -\sum_{E'}y_i(6p_i+f_i)+O(k^{-2}).
\]
The boundary identity for the action, the change in the prefactor, and the cofactor give
\[
 \begin{aligned}
 \epsilon_c:={}&L(B+c)-L(B)-\log p_c\\
 ={}&-6p_c-f_c+6\sum_Fp_i^2-6\sum_{E'}p_iy_i-\sum_{E'}y_if_i\\
 &+\sum_Fp_if_i^I-\sum_{j\in J}(\Gamma_{jj}-1)-1/k+O(k^{-2}).
 \end{aligned}
\]
The final identity of Section~\ref{sec:B4T08}, before approximation, implies
\[
 \sum_Fp_i(f_i-f_i^I)
 =-\sum_{j\in J}(\Gamma_{jj}-1)-1/k-p_cf_c.
\]
Since $p_cf_c=O(k^{-2})$, we obtain $\epsilon_c-\delta_c=O(k^{-2})$.
All error sums use $\sum_Ep=k$, $\sum_Fp=O(1)$, and Section~\ref{sec:B5T04}.
This step retains and cancels the terms of order $k^{-1}$.

The row conservation identity gives exactly
\[
 \sum_{c\in r}p_c\delta_c=0,\quad
 w_c=p_c(1+\delta_c)>0,\quad\sum_{c\in r}w_c=1.
\]
Section~\ref{sec:B5T02} implies $\delta_c=O(k^{-1})$. Hence
\[
 r_c=\epsilon_c-\log(1+\delta_c)=O(k^{-2}),\quad
 \rho=\mathbb E_w e^{r_c},\quad|\log\rho|\le\sup_c|r_c|.
\]
This proves the deterministic extension under the hypotheses of the theorem.
\end{proof}

\paragraph{Scope.}
The case of masked rooks is included through the representation with two
single-square diagonals per square in Section~\ref{sec:B5T01}. A child of an interior parent inherits
membership in an outer class with a margin. We make no claim that a trajectory remains in that class. We
exhibit here no residual
of queens in this class with $k=o(n)$; Corollary~\ref{cor:completo} only shows that the macroscopic
class contains the full board.
